\section{Orbital calibration and sextic consistency tests}
\label{app:orbital}

\subsection{Leading quartic information from axial actions}

Both axial period-two families satisfy
\[
\mcJ_H
=
-\frac{\delta^2}{16A}
+\bigO(\delta^3).
\]
With
\[
\bar{\mcJ}
=
\frac{\mcJ_Q+\mcJ_P}{2},
\]
define
\begin{equation}
\boxed{
A_{\rm orb}(\delta)
=
-\frac{\delta^2}{16\bar{\mcJ}}.
}
\label{eq:E-Aorb}
\end{equation}
Then
\[
A_{\rm orb}(\delta)=A+\bigO(\delta).
\]
This relation uses only the exact reduced symplectic actions of the axial periodic orbits.

\subsection{Leading quartic angular coefficient from physical Floquet exponents}

For the physical second return $F=P^2$,
\[
\kappa_{F,Q}=2\rho\delta+\bigO(\delta^2),
\qquad
\kappa_{F,P}=2\rho\delta+\bigO(\delta^2).
\]
Hence
\begin{equation}
\boxed{
\rho_{\rm orb}(\delta)
=
\frac{\kappa_{F,Q}+\kappa_{F,P}}{4\delta}
=
\rho+\bigO(\delta).
}
\label{eq:E-rhoorb}
\end{equation}
The corresponding orbital near-degeneracy parameter is
\[
\boxed{
\sigma_{\rm orb}=4-\rho_{\rm orb}^2.
}
\]
Once $A_{\rm orb}$ and $\rho_{\rm orb}$ are known, the quartic mixed coefficient can also be reconstructed as
\[
\boxed{
B_{\rm orb}
=A_{\rm orb}\left(2-\rho_{\rm orb}^2\right).
}
\]

\subsection{Richardson extrapolation}

If
\[
R(\delta)=R_\star+c_1\delta+\bigO(\delta^2),
\]
then the first Richardson transform
\begin{equation}
\boxed{
R^{[1]}(\delta)
=2R(\delta/2)-R(\delta)
}
\label{eq:E-richardson}
\end{equation}
removes the leading linear correction.  We apply \eqref{eq:E-richardson} to $A_{\rm orb}$ and $\rho_{\rm orb}$ on two interlaced dyadic ladders.  The convergence displayed in \cref{tab:orbital-quartic} provides an orbital validation of the quartic near-degeneracy.

\subsection{Why the subleading axial channel is not a critical-sextic invariant}
\label{app:axial-contamination}

A crucial point is that the detuned normal form is not exhausted by
\[
-\delta I+H_4+H_6.
\]
At the order relevant to the first subleading axial corrections one must retain
\begin{equation}
\boxed{
K_\delta
=
-\delta I
+H_4
+H_6
+\delta H_{4,1}
+\gamma\delta^2 I
+\bigO(\delta^3I,\delta^2I^2,\delta I^3,I^4),
}
\label{eq:E-extended-detuning}
\end{equation}
where $H_{4,1}$ is a quartic polynomial describing the first detuning derivative of the quartic germ in the chosen parameter-dependent normalization.

On an axial branch,
\[
I=\bigO(\delta).
\]
Therefore
\[
H_6=\bigO(\delta^3),
\qquad
\delta H_{4,1}=\bigO(\delta^3),
\qquad
\gamma\delta^2I=\bigO(\delta^3).
\]
The cubic action coefficient is thus not determined by $H_6$ alone.  The same contamination occurs in the $\delta^2$ correction to the physical Floquet exponent, because the Hessians of $\delta H_{4,1}$ and $\gamma\delta^2I$ contribute at precisely that order.

Consequently, formulas that attempt to infer the critical coefficients $d,e,f,g$ from the subleading axial action and Floquet coefficients while omitting $H_{4,1}$ and the quadratic detuning term are not valid in general.  The exact-return campaign detects this directly: the leading limits \eqref{eq:E-Aorb}--\eqref{eq:E-rhoorb} converge as expected, whereas the subleading coefficients are not described by the frozen critical sextic tensor alone.

This does not invalidate the critical Hamiltonian $H_6$.  It only means that a subleading \emph{detuned} axial reconstruction requires the parameter-dependent jet through the same weighted order.  We therefore do not use the axial subleading channel to estimate $G$ or the sextic anisotropy in the final paper.

\subsection{Elliptic action reconstruction of the diagonal sextic coefficient}

The near-diagonal scalar reduction is
\[
K_E(I)
=-\delta I+aI^2+GI^3,
\qquad
a=A\sigma<0.
\]
At an equilibrium,
\[
-\delta+2aI+3GI^2=0,
\]
while the reduced action is
\[
\mcJ_E=-aI^2-2GI^3.
\]
Eliminating $G$ gives
\[
3\mcJ_E=aI^2-2\delta I.
\]
The physical positive solution is
\begin{equation}
\boxed{
I
=
\frac{\delta-\sqrt{\delta^2+3a\mcJ_E}}{a},
}
\label{eq:E-I-from-action}
\end{equation}
and therefore
\begin{equation}
\boxed{
G_{\rm ell}
=
\frac{\delta-2aI}{3I^2}.
}
\label{eq:E-Gell}
\end{equation}

Unlike the contaminated axial subleading coefficients, this reconstruction is tested in the distinguished near-diagonal regime where
\[
I=\bigO(\delta^{1/2})
\]
in the outer sextic scaling.  There
\[
H_6=\bigO(\delta^{3/2}),
\]
whereas
\[
\delta H_{4,1}=\bigO(\delta^2),
\]
so the sextic term is asymptotically promoted ahead of the first detuning-dependent quartic correction.  This ordering is precisely what makes the near-diagonal branch the cleaner dynamical probe of $G$.

The exact-return values of $G_{\rm ell}$ are listed in \cref{tab:Gell}.  In the inner crossover region they agree with the frozen critical-jet value at the $10^{-3}$ level or better, with the expected gradual drift at larger detuning.

\subsection{Residual action and Floquet consistency}

At $\delta=0$ the residual exact cycle supplies an additional, independent sixth-order check.  Its reduced action is
\[
\mcJ_0^{\rm exact}
=-1.69102354131\times10^{-13},
\]
while the scalar frozen-coefficient prediction is
\[
\mcJ_0^{\rm scalar}
=-1.69009171482\times10^{-13}.
\]
The physical second-return angle is
\[
\vartheta_F^{\rm exact}
=1.62684506536\times10^{-4},
\]
compared with
\[
\vartheta_0^{\rm scalar}
=1.62689147576\times10^{-4}.
\]
Thus both invariant observables lie extremely close to the frozen sixth-order crossover endpoint.

\subsection{Logical hierarchy after the numerical audit}

The orbital information is therefore used in the following hierarchy:
\[
\boxed{
\text{leading axial action}\longrightarrow A,
}
\]
\[
\boxed{
\text{leading axial Floquet exponent}\longrightarrow\rho,
}
\]
\[
\boxed{
\text{near-diagonal action crossover}\longrightarrow G,
}
\]
while
\[
\boxed{
\text{subleading axial coefficients}
}
\]
are treated as mixed information involving both the critical sextic jet and detuning-dependent normal-form coefficients.  This separation is the convention used in the final numerical analysis.
